\section{Discussion}
Proving is now feasible but not yet deployable, and the gap is a \emph{dual obstruction}. The proof that is both feasible and post-quantum, the plain \emph{receipt} the zkVM emits, does not hide the holder's data (the first obstruction). The wrap SP1 adds to hide it does fit the $24\,\mathrm{GB}$ budget, but it relies on elliptic-curve cryptography, which a quantum computer breaks, so it is not post-quantum (the second obstruction). A future quantum adversary could later unmask the holder of a stored proof. No proof SP1's tools produce both hides its data and resists a quantum computer. The gap lies in the wraps SP1 ships, not in zkVMs as such. For the credential BDEC~\cite{10.1007/978-981-96-0957-4_3} built on PLUM, we prove a sharper limit. This route cannot give \emph{everlasting} anonymity (safety even against an attacker who stores a proof today and breaks it later with unlimited power), unless the proof first conceals its data with an extra step the toolkit does not provide~\cite{cgh2004rom}. The positive security results are conditional, resting on five named open assumptions together with standard ones inherited from the components.
